\subsection{Dependencias de la demostración}\label{sec:dependencias}
La tabla~\ref{tab:dependencias} indica en qué orden se pueden verificar las estimaciones.
En particular, el cambio de ley que transfiere la cota de llegada utiliza una estimación
obtenida previamente bajo la ley adaptativa.
\begin{center}
\small
\begin{tabular}{@{}L{.30\linewidth}L{.63\linewidth}@{}}
\toprule Resultado & Entradas y uso posterior\\ \midrule
Consistencia y permanencia &
Los resultados de~\cite{pc,tc} de la sección~\ref{sec:importados} fijan las clases, las tasas
de salida y el tamaño de los defectos de fila.\\
Llegada bajo \(\nu\) &
La positividad (Lema~\ref{lem:positividad-llegada}) permite tomar logaritmos.
Menores, coeficientes por soporte y cola (proposiciones~\ref{prop:menores-llegada},
\ref{prop:soportes-llegada} y~\ref{prop:norma-llegada}) se combinan con la geometría de los
tipos para estimar la media bajo \(\nu\).\\
Tasa y transferencia &
El Lema~\ref{lem:previa} aplica la identidad detenida a esa media a un único corte.
Con la tasa ya disponible, la densidad detenida lleva la cota a \(\mu\)
(Lema~\ref{lem:llegada-integrada}).\\
Interior y frontera &
En el interior se integra la suma de las dieciséis contracciones
(apéndice~\ref{app:interior}). En la frontera se separan los menores dobles,
la variación de un par y la interacción del cuadrado
(apéndices~\ref{app:boundary}, \ref{app:weights} y~\ref{app:transfer}).\\
Límite y coeficiente &
La identidad~\eqref{eq:signed-mean} reúne ambas medias.
La Proposición~\ref{prop:cortes-uniformes} controla sus errores; el
Lema~\ref{lem:dos-cortes} da el límite. El análisis del extremo identifica la
integral del Teorema~\ref{teo:principal}.\\
\bottomrule
\end{tabular}
\refstepcounter{table}\label{tab:dependencias}\par\smallskip
\textbf{Tabla \thetable.} Orden de los resultados y de los cambios de ley.
\end{center}
